% Generated draft: requires mathematical and source review.
\documentclass[11pt]{article}
\usepackage[margin=25mm]{geometry}
\usepackage{fontspec}
\setmainfont{Latin Modern Roman}
\usepackage{amsmath,amssymb,mathtools,unicode-math}
\setmathfont{Latin Modern Math}
\usepackage{enumitem,needspace,xurl,hyperref}
\setlength{\parindent}{0pt}
\setlength{\parskip}{5pt}
\setlength{\emergencystretch}{3em}
\newcommand{\N}{\mathbb{N}}
\newcommand{\Z}{\mathbb{Z}}
\newcommand{\Q}{\mathbb{Q}}
\newcommand{\R}{\mathbb{R}}
\newcommand{\C}{\mathbb{C}}
\newcommand{\F}{\mathbb{F}}
\newcommand{\A}{\mathbb{A}}
\newcommand{\PP}{\mathbb{P}}
\newcommand{\E}{\mathbb{E}}
\newcommand{\eps}{\varepsilon}
\newcommand{\dd}{\,\mathrm{d}}
\begin{document}
\section*{UnsolvedMath ID: 2305076}
\textbf{Source problem code:} catalogue entry
\subsection{Definitions and mathematical statement}
Let \(\mathbb D=\{z:|z|<1\}\) and \(H^\infty(\mathbb D)=\{g\text{ holomorphic on }\mathbb D:\sup_{\mathbb D}|g|<\infty\}\). Let \(\gamma\) be an arc with one endpoint \(1\), otherwise contained in \(\mathbb D\). The natural rotation convention is \(\gamma_\theta=e^{i\theta}\gamma\), an arc ending at \(e^{i\theta}\). With this convention, the requested property is
\[
\exists g\in H^\infty(\mathbb D)\quad\forall\theta\in[0,2\pi):\qquad
\lim_{z\to e^{i\theta},\ z\in\gamma_\theta}g(z)\text{ does not exist}.
\]
The source instead prints \(\gamma_\theta=e^{i\theta\gamma}\). This exponential-of-an-arc expression is not the same as rotation and need not be an arc inside the disk; the rotation formula above is an explicit possible correction, not a silent transcription change.

There is a second substantive issue: the source calls \(\gamma\) nontangential. In the usual Stolz-cone sense this means that, near \(1\),
\[
|1-z|\le C(1-|z|)\qquad(z\in\gamma)
\]
for some finite \(C\). Under the rotation reading, every rotated arc is then a nontangential approach at its endpoint. The usual almost-everywhere nontangential-limit theorem for \(H^\infty\) is incompatible with failure at \emph{every} endpoint in that interpretation. Thus the intended arc hypothesis or approach convention must be checked; no alternate class is guessed.

The source also states that the tangential-arc case has an affirmative answer, citing Collingwood--Lohwater. This comparison and its precise tangential convention are retained as source claims. The supplied OPEN-TRIAGE summary reports open status in the 2018 edition and no recent resolution located, but does not address the formula or nontangential-limit inconsistency.
\subsection{Short English statement}
Can a bounded analytic function fail to have a limit along every rotated copy of a prescribed approach arc? The printed rotation formula is malformed, and the standard nontangential reading conflicts with almost-everywhere boundary limits.
\subsection{Sources}
\begin{itemize}
\item Local UnsolvedMath registry, permanent ID \(2305076\), exact source code \texttt{catalogue entry}, ``Research Problems in Function Theory --- Problem 5.76.''
\item Hayman, \emph{Research Problems in Function Theory} (2018), Problem 5.76; attributed to W. Rudin and communicated by K. F. Barth. \url{https://arxiv.org/abs/1809.07200}. Source-TeX extraction; recorded access 2026-07-29.
\item The source mentions Rudin and Collingwood--Lohwater without complete references. The supplied OPEN-TRIAGE summary gives no separate current-status citation.
\end{itemize}
\end{document}
